\section{Experimental setup}

\subsection{Platform and task}

All experiments use a six degree of freedom UR3e manipulator model. The Cartesian task is a heart shaped trajectory in the $YZ$ plane with scale factor $0.008$. The trajectory is translated so that its initial point coincides with the end effector position generated by
\begin{equation}
  \q_0 =
  \begin{bmatrix}
    0 & -\pi/2 & \pi/2 & 0 & \pi/2 & 0
  \end{bmatrix}^{\top},
\end{equation}
which removes artificial startup mismatch in the offline studies.

Two execution modes are used at the current stage of the study.
\begin{itemize}[leftmargin=1.5em]
  \item \textbf{Offline mode:} numerical replay over a full \SI{20}{\second} cycle with the same kinematic model in controller and plant.
  \item \textbf{Real time control loop validation:} synchronous CoppeliaSim execution through the remote API. The main real time comparison uses \SI{10}{\second} cycles. An additional \SI{20}{\second} run of the original \methodthree{} is available as an auxiliary long period reference.
\end{itemize}

\subsection{Compared methods}

The compared methods are organized into three groups.
\begin{itemize}[leftmargin=1.5em]
  \item \textbf{Original baselines:} \methodone{}, \methodtwo{}, and \methodthree{}.
  \item \textbf{Formulation preserving solver transfers:} \methodoneDL{}, the default \methodtwoDL{}, and the tuned proposed \methodtwoDL{}.
  \item \textbf{Auxiliary structure matching studies:} \methodthreelcc{} and \methodthreepcg{}.
\end{itemize}

The purpose of the first group is to identify the baseline control structure. The second group tests whether the classic formulation can be improved through sampled DLCCZNN design. The third group explains why good raw numerical performance does not automatically make a method the main contribution of this paper.

\subsection{Parameter settings}

The original baseline settings are kept exactly as implemented in the integrated scripts. They are summarized in Table~\ref{tab:baseline_parameters}.

\begin{table}[H]
  \centering
  \caption{Baseline parameter settings used in the original integrated scripts.}
  \label{tab:baseline_parameters}
  \begin{threeparttable}
    \begin{tabular}{p{2.2cm}p{2.0cm}p{9.1cm}}
      \toprule
      Method & Step size $\tau$ & Main settings \\
      \midrule
      \methodone &
      $0.005$ &
      $k_{\mathrm{p}}=20$, $\mu=0.1$, $\rho=0.9$, $c_{\mathrm{clip}}=2.0$, $\gamma=80$. \\
      \methodtwo &
      $0.005$ &
      $k_{\mathrm{p}}=160$, $\mu=10$, $\rho=0.9$, $c_{\mathrm{clip}}=4.0$, $\gamma=1280$, $20$ inner substeps. \\
      \methodthree &
      $0.005$ &
      $k_{\mathrm{p}}=198$, $\alpha=10^{4}$, $\beta=10^{-3}$, $\varepsilon=10^{-9}$. \\
      \bottomrule
    \end{tabular}
  \end{threeparttable}
\end{table}

The main advanced settings used in this paper are listed in Table~\ref{tab:advanced_parameters}. \methodoneDL{} and \methodthreelcc{} are included with their implemented defaults and are treated only as auxiliary evidence, so their full parameter scans are omitted for brevity.

\begin{table}[H]
  \centering
  \caption{Main advanced settings for the proposed and auxiliary solver variants.}
  \label{tab:advanced_parameters}
  \begin{threeparttable}
    \begin{tabular}{p{3.0cm}p{8.3cm}p{2.4cm}}
      \toprule
      Method & Main settings & Role \\
      \midrule
      \methodtwoDL{} default &
      $k_{\mathrm{p}}=160$, $\mu_{\mathrm{a}}=10$, $\rho_{\mathrm{a}}=0.90$, $N_s=60$, $\gamma_{\mathrm{s}}=2560$ &
      direct transfer \\
      Proposed \methodtwoDL{} &
      $k_{\mathrm{p}}=220$, $\mu_{\mathrm{a}}=20$, $\rho_{\mathrm{a}}=0.85$, $N_s=80$, $\gamma_{\mathrm{s}}=4608$ &
      main method \\
      \methodthreepcg{} &
      $k_{\mathrm{p}}=198$, $w_{\mathrm{p}}=1.5\times10^{4}$, $w_{\mathrm{d}}=5\times10^{-4}$, $\varepsilon=10^{-10}$, $k_{\max}=6$ &
      auxiliary reference \\
      \bottomrule
    \end{tabular}
    \begin{tablenotes}[flushleft]
      \item The proposed \methodtwoDL{} setting is the real time drift focused recommendation obtained from the completed parameter study.
    \end{tablenotes}
  \end{threeparttable}
\end{table}

\subsection{Tuning protocol for the proposed method}

The retained \methodtwoDL{} setting was not selected by one pass trial and error. A two stage tuning procedure was carried out. First, a broad offline scan and a refined offline scan were executed to identify stable parameter families. This stage covered $81$ broad cases and $243$ refined cases. Second, representative cases from the stable family were replayed in real time CoppeliaSim control loops. This second stage covered $9$ real time cases and was used to choose the final retained setting.

The final selection principle is important. The retained configuration was chosen by prioritizing cycle end drift recovery under acceptable real time tracking error, rather than by selecting the smallest offline mean position error. This choice reflects the true control objective of drift free repetitive motion and will be revisited explicitly in the results discussion.

\subsection{Evaluation metrics}

For each run, the main reported metrics are the mean position error
\begin{equation}
  \frac{1}{N}\sum_{k=1}^{N}\|\x(\q_k)-\xd(k)\|_2,
\end{equation}
the final position error $\|\x(\q_N)-\xd(N)\|_2$, and the final drift norm $\|\q_N-\q_0\|_2$. The transient maximum absolute joint deviation is also monitored to distinguish normal motion excursion from true cycle end drift.

\subsection{Current status and mandatory physical validation}

The current manuscript reports offline data and simulator based real time control loop validation only. This is sufficient to establish the comparative logic of the paper, but it is not sufficient for a final Control Engineering Practice submission. A matching UR3e hardware protocol has already been prepared and must be executed before submission. The planned hardware study will use the same initial posture, the same heart trajectory, and the same cycle durations, and will report the same metrics together with the six joint angle recovery table at the end of each cycle. Additional hardware specific observations will include controller communication latency, repeated cycle accumulation, and the consistency between internal command trajectories and measured joint feedback.
